\documentclass[a4paper,12pt,twoside]{memoir}

% Castellano
\usepackage[spanish,es-tabla]{babel}
\selectlanguage{spanish}
\usepackage[utf8]{inputenc}
\usepackage[T1]{fontenc}
\usepackage{lmodern} % Scalable font
\usepackage{microtype}
\usepackage{placeins}
\usepackage{pdflscape}

\RequirePackage{booktabs}
\RequirePackage[table]{xcolor}
\RequirePackage{xtab}
\RequirePackage{multirow}

% Links
\PassOptionsToPackage{hyphens}{url}\usepackage[colorlinks]{hyperref}
\hypersetup{
	allcolors = {red}
}

% Ecuaciones
\usepackage{amsmath}

% Rutas de fichero / paquete
\newcommand{\ruta}[1]{{\sffamily #1}}

% Párrafos
\nonzeroparskip

% Huérfanas y viudas
\widowpenalty100000
\clubpenalty100000

% Imagenes
\usepackage{graphicx}
\newcommand{\imagen}[2]{
	\begin{figure}[!h]
		\centering
		\includegraphics[width=0.9\textwidth]{#1}
		\caption{#2}\label{fig:#1}
	\end{figure}
	\FloatBarrier
}

\newcommand{\imagenflotante}[2]{
	\begin{figure}%[!h]
		\centering
		\includegraphics[width=0.9\textwidth]{#1}
		\caption{#2}\label{fig:#1}
	\end{figure}
}



% El comando \figura nos permite insertar figuras comodamente, y utilizando
% siempre el mismo formato. Los parametros son:
% 1 -> Porcentaje del ancho de página que ocupará la figura (de 0 a 1)
% 2 --> Fichero de la imagen
% 3 --> Texto a pie de imagen
% 4 --> Etiqueta (label) para referencias
% 5 --> Opciones que queramos pasarle al \includegraphics
% 6 --> Opciones de posicionamiento a pasarle a \begin{figure}
\newcommand{\figuraConPosicion}[6]{%
  \setlength{\anchoFloat}{#1\textwidth}%
  \addtolength{\anchoFloat}{-4\fboxsep}%
  \setlength{\anchoFigura}{\anchoFloat}%
  \begin{figure}[#6]
    \begin{center}%
      \Ovalbox{%
        \begin{minipage}{\anchoFloat}%
          \begin{center}%
            \includegraphics[width=\anchoFigura,#5]{#2}%
            \caption{#3}%
            \label{#4}%
          \end{center}%
        \end{minipage}
      }%
    \end{center}%
  \end{figure}%
}

%
% Comando para incluir imágenes en formato apaisado (sin marco).
\newcommand{\figuraApaisadaSinMarco}[5]{%
  \begin{figure}%
    \begin{center}%
    \includegraphics[angle=90,height=#1\textheight,#5]{#2}%
    \caption{#3}%
    \label{#4}%
    \end{center}%
  \end{figure}%
}
% Para las tablas
\newcommand{\otoprule}{\midrule [\heavyrulewidth]}
%
% Nuevo comando para tablas pequeñas (menos de una página).
\newcommand{\tablaSmall}[5]{%
 \begin{table}
  \begin{center}
   \rowcolors {2}{gray!35}{}
   \begin{tabular}{#2}
    \toprule
    #4
    \otoprule
    #5
    \bottomrule
   \end{tabular}
   \caption{#1}
   \label{tabla:#3}
  \end{center}
 \end{table}
}

%
% Nuevo comando para tablas pequeñas (menos de una página).
\newcommand{\tablaSmallSinColores}[5]{%
 \begin{table}[H]
  \begin{center}
   \begin{tabular}{#2}
    \toprule
    #4
    \otoprule
    #5
    \bottomrule
   \end{tabular}
   \caption{#1}
   \label{tabla:#3}
  \end{center}
 \end{table}
}

\newcommand{\tablaApaisadaSmall}[5]{%
\begin{landscape}
  \begin{table}
   \begin{center}
    \rowcolors {2}{gray!35}{}
    \begin{tabular}{#2}
     \toprule
     #4
     \otoprule
     #5
     \bottomrule
    \end{tabular}
    \caption{#1}
    \label{tabla:#3}
   \end{center}
  \end{table}
\end{landscape}
}

%
% Nuevo comando para tablas grandes con cabecera y filas alternas coloreadas en gris.
\newcommand{\tabla}[6]{%
  \begin{center}
    \tablefirsthead{
      \toprule
      #5
      \otoprule
    }
    \tablehead{
      \multicolumn{#3}{l}{\small\sl continúa desde la página anterior}\\
      \toprule
      #5
      \otoprule
    }
    \tabletail{
      \hline
      \multicolumn{#3}{r}{\small\sl continúa en la página siguiente}\\
    }
    \tablelasttail{
      \hline
    }
    \bottomcaption{#1}
    \rowcolors {2}{gray!35}{}
    \begin{xtabular}{#2}
      #6
      \bottomrule
    \end{xtabular}
    \label{tabla:#4}
  \end{center}
}

%
% Nuevo comando para tablas grandes con cabecera.
\newcommand{\tablaSinColores}[6]{%
  \begin{center}
    \tablefirsthead{
      \toprule
      #5
      \otoprule
    }
    \tablehead{
      \multicolumn{#3}{l}{\small\sl continúa desde la página anterior}\\
      \toprule
      #5
      \otoprule
    }
    \tabletail{
      \hline
      \multicolumn{#3}{r}{\small\sl continúa en la página siguiente}\\
    }
    \tablelasttail{
      \hline
    }
    \bottomcaption{#1}
    \begin{xtabular}{#2}
      #6
      \bottomrule
    \end{xtabular}
    \label{tabla:#4}
  \end{center}
}

%
% Nuevo comando para tablas grandes sin cabecera.
\newcommand{\tablaSinCabecera}[5]{%
  \begin{center}
    \tablefirsthead{
      \toprule
    }
    \tablehead{
      \multicolumn{#3}{l}{\small\sl continúa desde la página anterior}\\
      \hline
    }
    \tabletail{
      \hline
      \multicolumn{#3}{r}{\small\sl continúa en la página siguiente}\\
    }
    \tablelasttail{
      \hline
    }
    \bottomcaption{#1}
  \begin{xtabular}{#2}
    #5
   \bottomrule
  \end{xtabular}
  \label{tabla:#4}
  \end{center}
}



\definecolor{cgoLight}{HTML}{EEEEEE}
\definecolor{cgoExtralight}{HTML}{FFFFFF}

%
% Nuevo comando para tablas grandes sin cabecera.
\newcommand{\tablaSinCabeceraConBandas}[5]{%
  \begin{center}
    \tablefirsthead{
      \toprule
    }
    \tablehead{
      \multicolumn{#3}{l}{\small\sl continúa desde la página anterior}\\
      \hline
    }
    \tabletail{
      \hline
      \multicolumn{#3}{r}{\small\sl continúa en la página siguiente}\\
    }
    \tablelasttail{
      \hline
    }
    \bottomcaption{#1}
    \rowcolors[]{1}{cgoExtralight}{cgoLight}

  \begin{xtabular}{#2}
    #5
   \bottomrule
  \end{xtabular}
  \label{tabla:#4}
  \end{center}
}


















\graphicspath{ {./img/} }

% Capítulos
\chapterstyle{bianchi}
\newcommand{\capitulo}[2]{
	\setcounter{chapter}{#1}
	\setcounter{section}{0}
	\chapter*{#2}
	\addcontentsline{toc}{chapter}{#2}
	\markboth{#2}{#2}
}

% Apéndices
\renewcommand{\appendixname}{Apéndice}
\renewcommand*\cftappendixname{\appendixname}

\newcommand{\apendice}[1]{
	%\renewcommand{\thechapter}{A}
	\chapter{#1}
}

\renewcommand*\cftappendixname{\appendixname\ }

% Formato de portada
\makeatletter
\usepackage{xcolor}
\newcommand{\tutor}[1]{\def\@tutor{#1}}
\newcommand{\course}[1]{\def\@course{#1}}
\definecolor{cpardoBox}{HTML}{E6E6FF}
\def\maketitle{
  \null
  \thispagestyle{empty}
  % Cabecera ----------------
\noindent\includegraphics[width=\textwidth]{cabecera}\vspace{1cm}%
  \vfill
  % Título proyecto y escudo informática ----------------
  \colorbox{cpardoBox}{%
    \begin{minipage}{.8\textwidth}
      \vspace{.5cm}\Large
      \begin{center}
      \textbf{TFG del Grado en Ingeniería Informática}\vspace{.6cm}\\
      \textbf{\LARGE\@title{}}
      \end{center}
      \vspace{.2cm}
    \end{minipage}

  }%
  \hfill\begin{minipage}{.20\textwidth}
    \includegraphics[width=\textwidth]{escudoInfor}
  \end{minipage}
  \vfill
  % Datos de alumno, curso y tutores ------------------
  \begin{center}%
  {%
    \noindent\LARGE
    Presentado por \@author{}\\ 
    en Universidad de Burgos --- \@date{}\\
    Tutor: \@tutor{}\\
  }%
  \end{center}%
  \null
  \cleardoublepage
  }
\makeatother

\newcommand{\nombre}{Rodrigo Garcia Martin} %%% cambio de comando

% Datos de portada
\title{Asistente de voz para Moodle
UBUVoiceAssistant 2.0}
\author{\nombre}
\tutor{Raúl Marticorena Sánchez}
\date{Julio 2021}

%in the preamble
%--------------------------------
  \usepackage[
    backend=biber,
    style=numeric,
  ]{biblatex}

 \addbibresource{bibliografia.bib}
%--------------------------------

\begin{document}

\maketitle


\newpage\null\thispagestyle{empty}\newpage


%%%%%%%%%%%%%%%%%%%%%%%%%%%%%%%%%%%%%%%%%%%%%%%%%%%%%%%%%%%%%%%%%%%%%%%%%%%%%%%%%%%%%%%%
\thispagestyle{empty}


\noindent\includegraphics[width=\textwidth]{cabecera}\vspace{1cm}

\noindent D. Raúl Marticorena Sánchez, profesor del departamento de Ingeniería Informática, área de Lenguajes y Sistemas Informáticos.

\noindent Expone:

\noindent Que el alumno D. \nombre, con DNI 71709733B, ha realizado el Trabajo final de Grado en Ingeniería Informática titulado UBUVoiceAssistant 2.0. 

\noindent Y que dicho trabajo ha sido realizado por el alumno bajo la dirección del que suscribe, en virtud de lo cual se autoriza su presentación y defensa.

\begin{center} %\large
En Burgos, {\large \today}
\end{center}

\vfill\vfill\vfill

% Author and supervisor
%\begin{minipage}{0.45\textwidth}
%\begin{flushleft} %\large
%Vº. Bº. del Tutor:\\[2cm]
%D. nombre tutor
%\end{flushleft}
%\end{minipage}
%\hfill
%\begin{minipage}{0.45\textwidth}
%\begin{flushleft} %\large
%Vº. Bº. del co-tutor:\\[2cm]
%D. nombre co-tutor
%\end{flushleft}
%\end{minipage}
\hfill

\vfill

% para casos con solo un tutor comentar lo anterior
% y descomentar lo siguiente
Vº. Bº. del Tutor:\\[2cm]
D. Raúl Marticorena Sánchez


\newpage\null\thispagestyle{empty}\newpage




\frontmatter

% Abstract en castellano
\renewcommand*\abstractname{Resumen}
\begin{abstract}
Los asistentes de voz son una tecnología de reciente aparición que se puede ver cada vez en más aspectos de la vida cotidiana, como en nuestros móviles, algunas páginas web, e incluso en las cajas registradoras automáticas que hay en algunos supermercados. Nos permiten realizar una enorme cantidad de funciones, como buscar información, enviar mensajes o navegar por diferentes páginas web, intentando que el funcionamiento sea lo más natural posible para el usuario, como si estuviese interactuando con otra persona.

En el campo de la educación se han utilizado muchas tecnologías diferentes hasta la fecha, siendo Moodle una de las más importantes, ya que permite ser un punto de unión entre alumnos y profesores, donde poder consultar todos los apuntes de las asignaturas, subir las tareas, realizar exámenes, comunicarse mediante foros y mensajería instantánea, y muchas otras funciones. También es importante destacar la facilidad para crear aplicaciones de terceros.

Debido a esto, surgió la idea de crear un asistente de voz para Moodle, para intentar hacer la experiencia más agradable. UBU\-Voice\-Assistant está programado en Python, y usa el asistente de código abierto \textit{Mycroft}. Se parte de una versión inicial y se buscan incluir varias mejoras en diferentes partes de la aplicación. En concreto, la instalación y configuración inicial se han simplificado en gran medida, se han realizado varios cambios en la interfaz para conseguir un aspecto más moderno, se han corregido una serie de errores y se han añadido algunas funcionalidades nuevas, como la posibilidad de interactuar con los mensajes privados de Moodle.
\end{abstract}

\renewcommand*\abstractname{Descriptores}
\begin{abstract}
asistente de voz, Mycroft, TTS, STT, Moodle, skill, Python
\end{abstract}

\clearpage

% Abstract en inglés
\renewcommand*\abstractname{Abstract}
\begin{abstract}
Voice assistants are a recently emerging technology that each day is being incorporated into more aspects of everyday life, such as cellphones, websites or even in the automatic cash registers in some supermarkets. They allow us to perform a huge number of actions, like looking for information, sending text messages or browsing the web, while trying to be as natural to the user as talking to another person.

In education, there are a lot of different technologies, staying Moodle as one of the most important ones, because students and teachers can use it as a nexus. There you can get the notes for the different subjects, upload assignments, take tests, chat with other people using forums and instant messaging, and many other things. Is worth noting the ease of creating third-party apps.

Because of that, we had the idea of creating a voice assistant for Moodle, trying to get a better user experience on the platform. I coded \textit{UBUVoiceAssistant} in Python, and it uses the open-source assistant Mycroft. I started from an initial version and I’m looking to include several improvements in different parts of the application. The installation and initial configuration are now way easier than before, and I redesigned the interface to get a more modern appearance. I also fixed a bunch of errors and added some additional features, such as the interaction with private messaging in Moodle.
\end{abstract}

\renewcommand*\abstractname{Keywords}
\begin{abstract}
voice assistant, Mycroft, TTS, STT, Moodle, skill, Python
\end{abstract}

\clearpage

% Indices
\tableofcontents

\clearpage

\listoffigures

\clearpage

\listoftables
\clearpage

\mainmatter
\capitulo{1}{Introducción}

Los asistentes de voz son una tecnología nueva que no se ha usado apenas hasta el momento para la educación. Sin embargo, se ha visto que puede ser muy interesante usarlas, debido a que es mucho más rápido y cómodo buscar determinadas informaciones con un comando de voz en vez de navegar por diferentes menús que pueden ser bastante complejos.

Este proyecto se basa en mejorar la primera versión de un asistente de voz que permitía conectarte a una plataforma Moodle (como UBUVirtual). Con respecto a la anterior versión, se busca simplificar el proceso de instalación, incluir una funcionalidad que te permita consultar los mensajes privados, ampliar el sistema de ayuda y hacer que la conversación sea algo más natural, haciendo que el programa vaya almacenando el contexto.

En esta memoria se irán explicando los conceptos necesarios para entenderlo, así como las tecnologías usadas para realizar el proyecto, los problemas que se han encontrado a lo largo del desarrollo, las decisiones tomadas y las soluciones para intentar solventarlos.

\capitulo{2}{Objetivos del proyecto}

Este apartado explica de forma precisa y concisa cuales son los objetivos que se persiguen con la realización del proyecto. Se puede distinguir entre los objetivos marcados por los requisitos del software a construir y los objetivos de carácter técnico que plantea a la hora de llevar a la práctica el proyecto.

\section*{Objetivos generales}
\begin{itemize}
    \item Mejorar la instalación del proyecto.
    \item Facilitar la internacionalización.
    \item Añadir sistemas de log.
    \item Mejorar el estilo visual.
    \item Conseguir mayor robustez y estabilidad en la aplicación.
    \item Mejorar la experiencia de usuario con una interacción más natural.
\end{itemize}

\section*{Objetivos técnicos}
\begin{itemize}
    \item Creación de un script de instalación, que sea capaz de descargar e instalar todos los componentes necesarios, y de poder desinstalar el programa, o poder actualizar a una nueva versión.
    \item Trabajo con interfaces gráficas.
    \item Usar HTML embebido en la aplicación para hacer que sea más visual que con los widgets del sistema operativo.
    \item Incorporación de una biblioteca de internacionalización para incluir nuevos idiomas de manera más sencilla, y sin tener que modificar el código.
    \item Crear una skill de Mycroft que aproveche el contexto para recordar las anteriores interacciones.
    \item Aprovechar el fuzzy matching para poder entender lo que el usuario quiere decir aunque el reconocimiento de voz no entienda bien las palabras.
\end{itemize}

\capitulo{3}{Conceptos teóricos}

En este apartado se desarrollan algunos de los principales conceptos teóricos que son necesarios para comprender el proyecto, como las definiciones de los componentes usados, o su funcionamiento.

\section{Asistente de voz}
Los asistentes de voz son programas que permiten a un usuario interactuar con una máquina, intentando que la comunicación entre ellos sea lo más natural posible, como si estuviésemos hablando con otra persona. Habitualmente para esto se usa directamente la voz, aunque también se suelen poder introducir órdenes escritas que realizan las mismas funciones. El asistente procesa las órdenes y responde de una forma similar.

\section{Wake word}
Para poder distinguir si una persona está hablando a un asistente de voz o lo está haciendo por cualquier otro motivo, se suelen usar una serie de palabras que el usuario debe pronunciar para que el asistente comience a escuchar y procese la orden que se le da. En este asistente es configurable, pero se suele usar \textit{``Hey Mycroft''}

\section{Skill}
Las skills son programas o aplicaciones que están diseñadas para un asistente de voz. Normalmente cada una de las skills se encarga de una función diferente, por ejemplo, una podría encargarse de mostrarte los próximos eventos del calendario mientras que otra se encarga de leerte las últimas noticias. Cada skill está compuesta por varios elementos, como pueden ser las utterance, los intents, los dialogs, los prompts o el contexto, que se describen más adelante.

\subsection{Utterance}
Una utterance es la frase que dice el usuario que sirve para activar una skill concreta y que en algunos causará que el asistente realice una acción determinada. Un posible ejemplo de utterance sería: ``Dime los próximos eventos de Sistemas Distribuidos''

\subsection{Intent}
Un intent son unas palabras clave del utterance que permiten al asistente determinar cuál es la acción que el usuario quiere realizar. Una skill puede tener asociados varios intents. Varios de ellos pueden lanzar una misma acción, consiguiendo una interacción más natural, ya que el usuario puede pedir lo mismo de diferentes formas, por ejemplo: ``Abrir el calendario'' o ``Consultar los eventos''.

\subsection{Dialog}
Este término es específico de Mycroft, aunque muchos otros asistentes usan otras herramientas para desempeñar la misma función. Los dialogs son las frases con las que responde el asistente a las peticiones que hace el usuario. También es posible que en vez de un dialog, la respuesta sea completamente dinámica y no se usen. En el ejemplo anterior, el dialog podría ser ``Los próximos eventos son: Entrega de la práctica 2 para el viernes 30 de Abril''.

\subsection{Contexto}
El contexto es una herramienta que se usa cada vez más dentro de los asistentes de voz ya que sirve para guardar parte de la información que se ha intercambiado en las anteriores preguntas y hacer que en las próximas interacciones los resultados ofrecidos por el programa estén relacionados. Por ejemplo, podríamos decir primero ``Dime los foros de Sistemas Distribuidos'' y si luego decimos ``Crear un hilo en los foros de esa asignatura'', el programa asociará ``esa asignatura'' con ``Sistemas Distribuidos''.

\subsection{Prompt}
Los prompts son preguntas que el asistente hace al usuario si el asistente necesita más información para completar la orden del usuario. Por ejemplo, si queremos crear un hilo en los foros de una asignatura, el asistente podría repetir el texto que ha entendido, y finalmente, usando un prompt, preguntar al usuario si el texto que ha entendido es correcto o no.

\section{Fuzzy matching}
Debido a que los sistemas de reconocimiento de voz a día de hoy distan bastante de ser perfectos, en ocasiones no detectan correctamente las palabras que el usuario dijo. Esto suele ocurrir de manera más habitual en el caso de palabras que no están recogidas en un diccionario o son usadas de manera frecuente, como por ejemplo, los nombres propios. Si se tiene una lista de opciones entre las que el usuario tiene que elegir una, se puede usar fuzzy matching para estimar cuál es la palabra más parecida de entre las opciones a lo que el usuario quiso decir. En nuestra aplicación, se pueden obtener los nombres de los contactos del usuario, estimar el parecido de la frase identificada por el reconocimiento de voz con cada uno de ellos y elegir el más similar.
\capitulo{4}{Técnicas y herramientas}

\section{API REST}
Una API (Application Program Interface) es un conjunto de definiciones y protocolos usados para desarrollar e integrar diferentes aplicaciones entre sí. En una API REST se establecen una serie de posibles peticiones HTTP que se pueden usar para intercambiar datos o realizar operaciones entre los diferentes componentes del sistema final.

\section{Moodle}
Moodle es una plataforma educativa gratuita y de código abierta que permite poner de manera sencilla a profesores y estudiantes en contacto, y ofrece multitud de herramientas para impartir y gestionar cursos a distancia. Se puede descargar desde su página oficial e instalar en prácticamente cualquier máquina en cuestión de minutos, o bien usar una versión de prueba como la Mount Orange School.
Moodle se usa en un gran número de colegios y universidades en todas partes del mundo, aunque en muchas de ellas se han implementado funciones específicas que son necesarias para esa institución concreta u ofrecer un diseño alternativo de la plataforma, pero manteniendo las funcionalidades básicas y las ventajas que ofrece, como la API REST, entre otras.
Este es el caso de la Universidad de Burgos, cuya plataforma se llama UBUVirtual.
\subsection{WebServices}
Web Services es el nombre que da Moodle a su API REST. Estos dan muchas facilidades a la hora de interactuar con la plataforma desde aplicaciones de terceros debido a que son peticiones muy simples y concretas. Son una gran ventaja con respecto a otras plataformas que no tienen este tipo de tecnologías ya que en esos otros casos habría que recurrir a técnicas bastante más complejas.
Como en la mayoría de APIs, es necesario primero obtener un código que identifica la sesión del usuario, habitualmente conocido como token, mediante las APIs que nos ofrecen para identificarnos como si estuviésemos usando la aplicación oficial de Moodle. A partir de ese punto, enviando nuestro token en cada petición a la API podemos interactuar con el resto de la plataforma y realizar las diferentes acciones, como consultar el calendario, los foros o los mensajes.

\section{Mycroft}
Mycroft es un asistente de voz gratuito y de código abierto. Actualmente está disponible tanto para sistemas basados en Linux como para Android, Raspberry Pi y dos dispositivos fabricados por Mycroft, el Mark 1 y el Mark II.
Es una aplicación que se ejecuta directamente en la máquina cliente, a diferencia de otros como Google Assistant que se ejecutan completamente en el servidor. También es diferente con respecto a otros muchos asistentes ya que Mycroft es modular y permite activar, desactivar y sustituir componentes, dependiendo de las necesidades.

Además, una de las principales ventajas de ser código abierto es que se pueden ver los componentes que lo forman de manera más sencilla y poder separar responsabilidades en el caso de tener que reimplementar alguno de sus componentes. Por defecto trae implementados seis componentes, aunque sólo nos importan cuatro de ellos para este proyecto.
\subsection{Audio}
Este componente se encarga tanto de transformar la voz del usuario al texto que se usa como entrada (llamado speech-to-text o STT) como de transformar el texto de salida a voz(text-to-speech o TTS).
Para STT actualmente se usa por defecto el motor de análisis de voz de Google, aunque también es importante destacar que están trabajando en un motor completamente open-source en colaboración con Mozilla, llamado DeepSpeech. A pesar de usar estos por defecto, también se pueden usar otros como el IBM Watson o el servicio de wit.ai. Debido a que para usar la mayoría de estos motores se necesita un modelo entrenado previamente, que puede llegar a ser muy pesado y habitualmente es imposible descargar al ser software privativo, en la mayoría de estas opciones es necesario recurrir a un servidor en Internet que analice los clips de sonido que se han grabado en el cliente.
Con respecto a TTS, algunas de las configuraciones más básicas se pueden hacer desde su página web. Ahí te permiten seleccionar tres motores, Mimic 1 si seleccionas British Male, Mimic 2 si seleccionas American Male o Google Voice para usar la API de Google TTS. Se pueden configurar algunos otros como eSpeak, Microsoft Azure o Amazon Polly editando la configuración de forma local. Ya que algunos de estos motores son más ligeros e incluso hay opciones de código abierto, es posible generar y reproducir la respuesta sin necesidad de conexión a Internet.
\subsection{Voice}
Este módulo es el que se encarga de detectar la Wake Word dicha por el usuario. Por defecto se usa Precise, que es una red neuronal entrenada con sonidos y que permite usar cualquier frase especificando los fonemas usados. Sin embargo, en caso de que este motor no esté disponible, ya sea porque es incompatible o porque el equipo no tiene tanta potencia de cálculo, es posible usar PocketSphinx que es más ligero pero sólo funciona correctamente en inglés.
\subsection{Skills}
Esta parte de Mycroft es un servicio que se encarga de relacionar el utterance entendido por el módulo de audio con alguno de los intents definidos en las skills que haya instaladas. También se encarga de analizar alguna de enlazar alguna de las palabras de la utterance con las variables que hay en los intents (por ejemplo, en "Dime los próximos eventos de Sistemas Distribuidos", asigna "Sistemas Distribuidos" a la variable de curso). Esto se consigue gracias a un intent parser.
Actualmente hay dos, Adapt y Padatious. Padatious está basado en redes neuronales y es bastante sencillo de usar, ya que tiene una buena documentación, se pueden declarar los intents de manera muy sencilla y gran parte del trabajo lo hace la red neuronal de manera transparente al programador. Adapt, en cambio, es un sistema bastante más ligero que está pensado para dispositivos con menos capacidad de cómputo ya que la mayor parte del trabajo se hace en el momento de programar la skill, lo que junto a su mala documentación, hace que sea más complejo de usar. Sin embargo, el primero todavía no soporta contextos, posiblemente debido a que es un intent parser bastante más complejo y es necesario más tiempo de desarrollo.
\subsection{MessageBus}
El MessageBus es un módulo que permite que los componentes anteriormente mencionados se comuniquen entre sí. Se basa en un websocket que se usa para transferir mensajes entre las diferentes partes de Mycroft e incluso permite integrar aplicaciones de terceros. Por ejemplo, cuando el usuario dice una frase, esta se interpreta en el módulo de Audio y se transfiere el texto al módulo de Skills mediante el bus, y este le devuelve la respuesta a través del bus, que se transforma a voz y se reproduce.

\section{Qt5}
Qt es un framework orientado a objetos que permite desarrollar interfaces gráficas, que puede usarse tanto en Windows como en Mac o Linux, entre otros, sin tener que realizar grandes cambios en la aplicación. Gran parte del código está hecho y pensado para C++, aunque hay un gran número de bindings que permiten usar sus ventajas en otros lenguajes de programación. Este framework está disponible bajo dos tipos de licencias, una de código abierto y otra comercial.
\subsection{Qt Designer}
Qt Designer es una aplicación que permite desarrollar de manera muy rápida e intuitiva prototipos para interfaces que usen la biblioteca Qt. Permite usar un gran número de Widgets que trae la biblioteca, como campos de texto, imágenes, casillas, o desplegables. Sin embargo, para funcionalidades más avanzadas es necesario hacerlo mediante código. La aplicación permite convertir los prototipos a código fuente o exportarlo a un formato que se puede cargar en muchos de los lenguajes soportados.

\section{Python}
Python es un lenguaje de programación interpretado y con tipado dinámico. Es multiplataforma ya que se puede usar en Windows, Mac y Linux y también multiparadigma, porque se puede usar para hacer programación imperativa, orientada a objetos y en las últimas versiones también para programación funcional.
En los últimos años se ha vuelto extremadamente popular debido a que es un lenguaje sencillo de entender, tiene una gran cantidad de bibliotecas para realizar todo tipo de aplicaciones y además permite crear programas en un menor tiempo que otros lenguajes.
Algunas de las bibliotecas reseñables que he usado son las siguientes:
\subsection{PyQt5}
Esta biblioteca es un binding no oficial para usar la biblioteca de interfaces gráficas que hemos elegido, Qt. Sin embargo, pese a no ser oficial, es el más completo disponible y además se puede usar sin muchos problemas siguiendo la documentación oficial de Qt.
\subsection{requests}
Requests nos ofrece la posibilidad de hacer peticiones HTTP de manera más sencilla y más elegante que usando las bibliotecas por defecto de Python. Se ha usado para interactuar con la API de Moodle de una manera eficaz.
\subsection{gettext}
Gettext es una biblioteca que viene en el paquete por defecto de Python que sirve para hacer bastante más sencilla la internacionalización de un programa. Esta biblioteca permite cargar las frases de texto a partir de archivos codificados en un estándar usado habitualmente en sistemas GNU.
\subsection{babel}
Babel trae una serie de herramientas relacionadas con el idioma. En concreto, permite obtener de manera automática el nombre del idioma sin tener que añadirlo explícitamente en cada idioma nueva que añadamos al programa.
\subsection{fuzzywuzzy}
Esta biblioteca permite simplificar el proceso del fuzzy matching al tener que comparar el nombre dicho por el usuario con los elementos que encontramos en la plataforma, especialmente útil para tratar con nombres propios.

\section{Bash}
Bash es un lenguaje de órdenes usado ampliamente en los sistemas Linux. Este lenguaje permite al usuario escribir órdenes en forma de texto desde una terminal. Nos ofrece numerosas herramientas, siendo una de las más importantes la posibilidad de ejecutar comandos de manera automatizada desde un archivo de texto, conocido habitualmente como script. En el proyecto se ha usado para crear la herramienta de instalación.

\section{LaTeX}
LaTeX es un sistema de composición de textos, usado habitualmente para artículos y libros científicos ya que permite conseguir alta calidad tipográfica. Una de las principales características es que procesa archivos de texto plano en los que se incluyen marcas que indican el formato que tiene que usarse en el documento final.
Hay múltiples paquetes que permiten crear documentos usando este sistema, siendo Tex Live una de las más habituales y completas.

\section{Metodología de desarrollo}
Para desarrollar el proyecto se ha hecho uso de una de las metodologías ágiles más populares, SCRUM, aplicándole algunos cambios que son necesarios. A lo largo del proyecto se han ido realizando una serie de sprints, siendo los primeros de dos semanas y los finales de sólo una para arreglar los últimos errores. También se han sustituido las reuniones diarias por una reunión al principio y al final de cada sprint. Se han seguido también algunas prácticas como la integración continua, en la que se realizan pruebas del código de forma habitual a lo largo del desarrollo.

\section{Git}
Git es un programa de control de versiones ampliamente utilizado para gestionar el código fuente durante el desarrollo de aplicaciones. Este sistema se puede usar de forma local o también se puede sincronizar el contenido de los proyectos en un repositorio, permitiendo a múltiples personas colaborar. Una de las funcionalidades clave es el sistema de ramas, que permiten trabajar en diferentes funciones independientes entre sí al mismo tiempo.
\subsection{Github}
Github es una de los principales servidores que usan Git para el control de versiones. Esta plataforma aloja gran cantidad de repositorios de código abierto de forma gratuita. Además, permite crear Issues o tareas en las que ir organizando el trabajo a realizar. Ofrece una herramienta muy potente para integración continua llamada Github Actions. Además, Github tiene una API que permite a otros desarrolladores añadir funcionalidades extra para adaptarlo a tus necesidades.
\subsection{Zenhub}
Zenhub es una de las múltiples aplicaciones integradas con Github que existen para poder crear un tablero Kanban para organizar las tareas y usar otra serie de funcionalidades útiles relacionadas con SCRUM y otras metodologías ágiles, como la gestión de sprints o la generación de gráficas que permiten conocer la evolución del proyecto a lo largo del tiempo.

\section{Visual Studio Code}
VSCode es un entorno de desarrollo integrado, o IDE, que permite editar archivos de texto plano, como el código fuente, y además cuenta con una serie de herramientas que permiten hacer más cómodo y rápido el trabajo y detectar algunos errores antes de llegar a probar el código. Está desarrollado por Microsoft y tiene una gran comunidad que ha desarrollado numerosas extensiones que añaden algunas funcionalidades muy útiles. Las más relevantes son:
\subsection{Code Spell Checker}
Esta extensión es un corrector que analiza los ficheros que tengas abiertos y marca las palabras con errores tipográficos. Soporta varios idiomas, entre ellos el inglés y el español, ambos usados en el proyecto.
\subsection{GitLens}
GitLens añade una serie de funcionalidades adicionales a la potente integración con Git ya incluida en el IDE. En concreto, mejora la gestión de ramas y la visualización de cambios en los commits anteriores.
\subsection{LaTeX Workshop}
Esta extensión permite trabajar de forma sencilla con los documentos LaTeX, añadiendo algunas plantillas de autocompletado para las marcas de capítulo, sección o tablas, y añadiendo la posibilidad de generar los archivos PDF finales de la memoria una vez acabada la redacción.
\subsection{Pylance y Python}
Estas dos extensiones ofrecen funcionalidades imprescindibles para trabajar con Python dentro de Visual Studio Code, siendo algunas de las más importantes el autocompletado, las ayudas de tipos o la búsqueda de errores y debugging.
\subsection{Python Docstring Generator}
Esta última extensión permite generar parte de la documentación incluida en cada método a partir del código que hemos programado anteriormente. Es especialmente importante porque dejar una buena documentación mejora la calidad del código.

\section{Análisis del software}
Uno de los procesos que se debe hacer en el desarrollo de software es su análisis. En este proceso se comprueban algunos aspectos del mismo, como la calidad, mantenibilidad y seguridad. Debido a que es un proceso muy importante, han surgido multitud de herramientas, entre las cuales destacaré:
\subsection{SonarCloud}
Es un servicio gratuito en la nube que permite analizar y visualizar mediante una página web algunas de las características más importantes del análisis del software.
\subsection{Pylint}
Es una herramienta gratuita y de código abierto que se ejecuta de forma local e informa al desarrollador acerca de algunos posibles errores en el código que podrían pasar inadvertidos y podrían causar problemas graves en un futuro si no se corrigen.
\subsection{Mypy}
Es una herramienta de código abierto que analiza los archivos fuente para avisar de posibles incoherencias de tipos que podrían causar errores graves más adelante.
\capitulo{5}{Aspectos relevantes del desarrollo del proyecto}

En este apartado se recogen los aspectos más importantes del desarrollo del proyecto, como los problemas que han ido surgiendo y las diferentes decisiones que se han ido tomando.

\section{Uso de contenedores Docker}
En un primer momento se consideró la posibilidad de usar contenedores Docker tanto para la distribución de Mycroft como para la interfaz gráfica para facilitar el proceso de instalación.
Mycroft ofrece un contenedor ya creado con su aplicación, junto a las instrucciones de uso. Al instalarlo y ejecutarlo durante las primeras semanas del proyecto, todo funcionaba correctamente. Sin embargo, al probar a realizar la instalación en otra máquina, el contenedor no captaba sonido del micrófono y en ocasiones fallaba al reproducir las respuestas de Mycroft. Al reinstalar el contenedor de Mycroft en la máquina en la que anteriormente funcionaba correctamente, también dejó de captar el sonido y no vi manera de solucionarlo. Por esto, se optó por usar la aplicación nativa de Mycroft para Linux.
Con respecto a la interfaz gráfica, al ser más complicado lanzar de manera automatizada Mycroft, ubicado en el equipo host, desde un contenedor en el que se ubicase la interfaz, se desestimó completamente el uso de Docker.

\section{Creación del instalador}
Debido a la compleja instalación que había en la versión anterior del proyecto, y al no poder usar contenedores ya listos, se opta por crear un instalador que automatice la obtención e instalación de las dependencias, la copia de los archivos a las ubicaciones adecuadas y la creación de accesos directos a la aplicación. Para realizarlo se opta por crear un \textit{script} de línea de comandos para la primera versión, eligiendo Bash debido a que se usa habitualmente para esta finalidad y a que está disponible en la gran mayoría de sistemas. Debido a la diferente arquitectura y popularidad de las diferentes distribuciones Linux, se opta por Ubuntu como la principal distribución a soportar, aunque el funcionamiento será también correcto en otras que derivan de ella.

\section{Mejoras de la interfaz}
En la versión anterior del proyecto se creó una interfaz gráfica, pero los archivos que permitían editarla mediante QtDesigner no se subieron al repositorio. Al ser conveniente realizar algunas mejoras de estabilidad y de diseño, se opta por crear una segunda versión de la interfaz desde cero, pero manteniendo un diseño lo más simple posible. En la nueva interfaz se consigue que sea más estable que la anterior, evitando algunos de los cuelgues que antes se producían en algunas partes de la ejecución. También se crea un nuevo apartado de la interfaz que te guía en el proceso del emparejamiento con los servicios de Mycroft, paso necesario para poder usar el reconocimiento de voz. Finalmente, se opta por hacer un rediseño de la interfaz de chat, pasando de usar en algunos apartados los \textit{widgets} nativos que ofrece Qt a usar HTML embebido que ofrece mayor flexibilidad a la hora de crear los bocadillos de conversación.

\section{Localización}
Anteriormente se podía seleccionar entre inglés o español como idioma para usar la aplicación, que estaban incluidos en el código fuente de la aplicación. Debido a que esto es una mala práctica, se opta por usar una biblioteca de localización que cargue cada uno de los textos a mostrar desde archivos externos al código, lo que también hace más fácil añadir nuevos idiomas posteriormente sin tener que modificar nada de código. Se elige \textit{gettext} como biblioteca a usar debido a la existencia de software como Poedit que hace sencilla la generación de nuevos archivos de idioma.

\section{Uso de contexto}
También se busca mejorar la interacción con el asistente, y se considera empezar a trabajar con el contexto para que el asistente recuerde algunos elementos a lo largo de la conversación. En las anteriores skills que se habían realizado para el proyecto, se había usado Padatious como intent parser pero al revisar la documentación, todavía no es posible usar el contexto. Al ver este problema, se opta por realizar una nueva skill usando Adapt.
Si bien es cierto que ambos \textit{intent parsers} permiten acceden a un gran número de funciones comunes proporcionadas por Mycroft, la manera de crear \textit{skills} es completamente diferente. Destacar que es más complicado crearlas usando Adapt debido a que la documentación es considerablemente más escasa y que es imprescindible usar expresiones regulares para definir cada uno de los \textit{intents}, donde además hay que indicar también los posibles signos de puntuación.

\section{Instalación en Windows}
Previamente, uno de los principales puntos débiles de este proyecto era la necesidad de tener un sistema operativo Linux para poder usar el asistente. Sin embargo, en las últimas versiones de Windows 10 se ha incorporado una funcionalidad (Windows Subsystem for Linux, o WSL) con la que se puede ejecutar una versión ligera de Linux dentro de ese sistema operativo, pudiendo ejecutar nuestra aplicación sin necesidad de instalar una máquina virtual completa. Por el momento, es necesario instalar en Windows un servidor X11 para mostrar interfaces gráficas y Pulseaudio para grabar y reproducir sonido, pero Microsoft está trabajando para que en versiones futuras del sistema, esto ya venga incluido.
\capitulo{6}{Trabajos relacionados}

Si bien es cierto que este campo estaba prácticamente inexplorado hace unos pocos años, cada poco tiempo aparece algún proyecto que pone a prueba las capacidades de los asistentes en el campo de la educación. Algunos de los más interesantes son los siguientes:

En primer lugar tenemos los predecesores del proyecto actual, realizados por la Universidad de Burgos. Podemos encontrar tanto \href{https://github.com/adp1002/UBUVoiceAssistant}{la versión anterior de este mismo proyecto} como \href{https://github.com/cgc0045/TFG-UBUassistant}{una aplicación para Android} que permite hacer preguntas sobre la información de la página web de la universidad.

En segundo lugar podemos encontrar algunos \textit{chatbots} que permiten interactuar con Moodle de manera bastante similar al UBUVoiceAssistant. Entre los más destacables se encuentra \href{https://github.com/AnyTimeTraveler/moodlebot}{una integración con Telegram}, o \href{https://education.microsoft.com/en-us/resource/3dffb3a8}{el plugin para Microsoft Teams}, donde además del bot, encontramos la posibilidad de ver los cursos de manera visual en el cliente de mensajería. También podemos encontrar una \href{http://libres.uncg.edu/ir/asu/f/Melton_Michelle_2019_Thesis.pdf}{skill para Alexa}, que tiene un propósito y una funcionalidad parecida a nuestro proyecto.

También podemos encontrar una serie de asistentes que permiten interactuar con otras plataformas universitarias. Por ejemplo, tenemos \href{https://www.admithub.com/case-study/how-georgia-state-university-supports-every-student-with-personalized-text-messaging/}{Pounce} para interactuar con la \textit{Georgia State University}, o \href{https://www.youtube.com/watch?v=zsRPuU53E74}{Genie} en la \textit{Deakin University}. Es necesario mencionar el caso de \textit{Georgia Tech} \cite{georgiatech}, en el que, durante uno de sus cursos online se usaron dos \textit{chatbots} que sustituyeron a profesores de apoyo para responder algunas de las preguntas más habituales.

Finalmente, cabe destacar el proyecto de la \textit{Sant Louis University} \cite{alexauni}, donde instalaron un total de 2300 dispositivos \textit{Amazon Echo} con los que la comunidad universitaria puede hacer preguntas sobre la universidad a Alexa, sin necesidad de usar un dispositivo propio.

\begin{landscape}
Se puede ver una comparación de las características en la tabla \ref{tabla:CompCaracteristicas}

\tablaSmall{Comparación de características}{c c c c c c}{CompCaracteristicas}{
    Nombre del & Lenguaje de & Tipo de & Dispositivos & Complejidad de & Complejidad \\
    proyecto & programación & información & compatibles & instalación & de uso \\
}{
    UBUVoiceAssistant 2 & Python & Dinámica (Moodle) & Windows y Linux & Baja & Moderada\\
    UBUVoiceAssistant 1 & Python & Dinámica (Moodle) & Linux & Alta & Moderada \\
    UBUAssistant & Java & Estática & Android & Baja & Moderada \\
    Bot de Telegram & Java & Dinámica (Moodle) & PC y móvil & Alta & Alta \\
    MS Teams & PHP & Dinámica (Moodle) & PC y móvil & Alta & Baja \\
    Skill de Alexa & Java & Dinámica (Moodle) & Móvil & Moderada & Moderada \\
    Georgia State University & Desconocido & Dinámica & Móvil & Baja & Baja \\
    Deakin University & Desconocido & Estática & Móvil & Baja & Moderada \\
    Georgia Tech & Desconocido & Estática & Aula virtual & Ninguna & Baja\\
    Sant Louis University & Java & Estática & Amazon Echo/Dot & Ninguna & Moderada \\
}

\end{landscape}
\capitulo{7}{Conclusiones y Líneas de trabajo futuras}

\section{Conclusiones}


\section{Líneas de trabajo futuras}
En este apartado se incluyen algunas de las ideas que surgieron en la anterior versión del proyecto que no se han podido cumplir, junto a otras  nuevas que han ido surgiendo a lo largo del desarrollo.

Por el momento sólo es posible usar la aplicación en Linux o en Windows 10 mediante el WSL, pero no se puede usar en MacOS ni en móviles, donde quizá es más útil. Por el momento es necesario esperar a que avance el desarrollo de Mycroft para que sea viable usarlo en esas plataformas. También se podría optar por intentar conseguir una funcionalidad similar mediante otras tecnologías que sí sean multiplataforma, o crear una aplicación cliente ligero que se conecte a un servidor de la universidad donde se procesen todas las peticiones de los diferentes usuarios.

También podría ser interesante la posibilidad de trabajar sin conexión. Sería necesario descargar los datos obtenidos de Moodle a una serie de archivos o una base de datos local que se podrían consultar más adelante cuando no haya una conexión disponible.


%Where the bibliography will be printed
  \printbibliography

\end{document}
